\input{preamble}

\section*{Scattering part 4}

From the previous section we have
\begin{equation*}
f(\theta)=-\frac{2m\beta}{p^2+\epsilon^2\hbar^2}
\end{equation*}

where $V(r)$ is the Yukawa potential
\begin{equation*}
V(r)=\frac{\beta}{r}\exp(-\epsilon r)
\end{equation*}

For the Coulomb potential
\begin{equation*}
V(r)=-\frac{q_1q_2}{4\pi\varepsilon_0r}
=-\frac{Ze^2}{4\pi\varepsilon_0r}
=-\frac{Z\alpha\hbar c}{r}
\end{equation*}

we substitute $\epsilon=0$ and $\beta=-Z\alpha\hbar c$ to obtain
\begin{equation*}
f(\theta)=\frac{2mZ\alpha\hbar c}{p^2}
\end{equation*}

Substitute
\begin{equation*}
p^2=4mE(1-\cos\theta)
\end{equation*}

to obtain
\begin{equation*}
f(\theta)=\frac{Z\alpha\hbar c}{2E(1-\cos\theta)}
\end{equation*}

Hence the Rutherford cross section formula
\begin{equation*}
\frac{d\sigma}{d\Omega}=|f(\theta,\phi)|^2
=\left[\frac{Z\alpha\hbar c}{2E(1-\cos\theta)}\right]^2
\end{equation*}

Noting that
\begin{equation*}
(1-\cos\theta)^2=4\sin^4(\theta/2)
\end{equation*}

we can also write
\begin{equation*}
\frac{d\sigma}{d\Omega}=\left(\frac{Z\alpha\hbar c}{2E}\right)^2
\frac{1}{4\sin^4(\theta/2)}
\end{equation*}

\href{https://georgeweigt.github.io/examples/scattering-part-4-demo.html}{Eigenmath script}

\end{document}
